\textbf{1. Kiến thức cần nhớ}
\begin{itemize}
\item \textbf{Một số hệ thức về cạnh và góc trong tam giác vuông}:
Cho tam giác $ABC$ vuông tại $A$, cạnh huyền $a$, các cạnh góc vuông $b, c$:
\begin{itemize}
\item Mỗi cạnh góc vuông bằng \textbf{cạnh huyền nhân với sin góc đối} hoặc \textbf{nhân với côsin góc kề}:
$$\boxed{b = a \cdot \sin B = a \cdot \cos C};\quad \boxed{c = a \cdot \sin C = a \cdot \cos B}$$
\item Mỗi cạnh góc vuông bằng \textbf{cạnh góc vuông kia nhân với tan góc đối} hoặc \textbf{nhân với côtan góc kề}:
$$\boxed{b = c \cdot \tan B = c \cdot \cot C};\quad \boxed{c = b \cdot \tan C = b \cdot \cot B}$$
\end{itemize}
\item \textbf{Giải tam giác vuông}: Là tìm tất cả các cạnh và các góc chưa biết của một tam giác vuông khi đã biết trước ít nhất hai yếu tố (trong đó phải có ít nhất một cạnh).
\end{itemize}

\textbf{2. Phương pháp giải các dạng}

%% DANG 9H4?12-1
\textbf{Dạng 1. Giải tam giác vuông}
\begin{enumerate}
\item \textbf{Trường hợp 1 (Biết một cạnh và một góc nhọn)}:
\begin{itemize}
\item Tính góc nhọn còn lại: $\widehat{C} = 90^\circ - \widehat{B}$.
\item Dùng hệ thức lượng giác (sin, cos hoặc tan) để tính hai cạnh còn lại.
\end{itemize}
\item \textbf{Trường hợp 2 (Biết hai cạnh)}:
\begin{itemize}
\item Dùng định lí Pythagore để tính cạnh thứ ba.
\item Dùng tỉ số lượng giác (sin, cos hoặc tan) để tính một góc nhọn, rồi suy ra góc nhọn còn lại.
\end{itemize}
\end{enumerate}

%% DANG 9H4?12-2
\textbf{Dạng 2. Tính độ dài, số đo góc trong hình}
\begin{enumerate}
\item Kẻ thêm đường cao để tạo ra các tam giác vuông.
\item Áp dụng hệ thức cạnh và góc trong từng tam giác vuông để tính đoạn thẳng trung gian.
\item Suy ra độ dài đoạn thẳng hoặc góc cần tìm.
\end{enumerate}

%% DANG 9H4?12-3
\textbf{Dạng 3. Ứng dụng thực tế (đo chiều cao, khoảng cách)}
\begin{enumerate}
\item \textbf{Bài toán xác định chiều cao vật thể} (cột cờ, tháp, ngọn hải đăng, tòa nhà):
\begin{itemize}
\item Đặt giác kế cách chân tháp một khoảng $a$.
\item Đo góc nâng $\alpha$, tính chiều cao phần trên $h_1 = a \cdot \tan \alpha$.
\item Chiều cao toàn bộ: $H = h_1 + h_{\text{giác kế}}$.
\end{itemize}
\item \textbf{Bài toán góc nghiêng dốc, đường bay máy bay}:
\begin{itemize}
\item Chiều cao đạt được: $h = s \cdot \sin \alpha$ (với $s$ là quãng đường đi).
\item Quãng đường đi: $s = \dfrac{h}{\sin \alpha}$.
\end{itemize}
\item Chú ý đổi đồng nhất đơn vị đo và làm tròn số theo yêu cầu của đề bài.
\end{enumerate}
